\chapter{Ethical Implications and Social Impact}
\setcounter{section}{0}
\setlength{\parindent}{0pt}
\section{AI-Based Systems Ethical Impact}

The rapid advancements in artificial intelligence systems have raised a new debate
on its ethical implications. Traditionally, software programs execute explicit and
deterministic instruction, which are given and controlled by a human. Also, the same
input always gave the same output, but if a misbehaviour occurred, one could always
traceback to a specific block of code. However, agentic systems change the current
paradigm, they can interpret natural language, reason over context and plan a sequence
of actions accordingly. They can interact with the real world on behalf of the user
or fully autonomously, with the added disadvantage of making it very difficult the
entire thought process it had on why doing certain actions. These AI models are
starting to shift from tools to operators.

This shift applies particularly to desktop automation, where an agent can interact
with the graphical interface and access everything visible on screen, this means
sensitive data can be manipulated in the worse case scenario. Furthermore, the actions
taken might be irreversible, such as deleting from the recycling bin a file. Therefore,
a misinterpreted command does not only return a wrong answer but might expose
confidential data about the user. As seen the ethical stakes are considerably higher
than a regression model.

To address these concerns, several legal frameworks regulate the development and
use of AI systems of this kind.

\begin{itemize}
    \item \textbf{General Data Protection Regulation (GDPR, Regulation (EU) 2016/679)}
    \cite{gdpr}: The GDPR governs the processing of personal data in the EU. Its
    Article~25 introduces \textit{privacy by design and by default}, requiring data
    protection to be embedded into a system's architecture from the beginning.
    \vspace{2.5em}
    \item \textbf{Spanish Organic Law on Data Protection (LOPDGDD, Ley Orgánica
    3/2018)} \cite{lopdgdd}: The LOPDGDD is the national transposition of the GDPR
    into Spanish law. It reinforces the same principles of data minimisation, purpose
    limitation, and privacy by design, and is directly applicable given that Spain is
    the primary development and deployment context of this project.

    \item \textbf{European Artificial Intelligence Act} \cite{ai_act}: This law aims
    to regulate the development and use of AI systems, ensuring that they are safe,
    transparent, and respectful of fundamental rights. It classifies AI applications
    into different risk levels and establishes strict requirements for high-risk
    systems, with the purpose of guaranteeing that they are used in a fair and
    ethical manner.
\end{itemize}

\section{Ethical Dimensions Descriptions}
\label{sec:ethical_dimensions}

Google's AI Principles \cite{google_principles} are structured around seven concrete
criteria that map directly onto the design decisions made throughout this project.
Their focus on privacy, safety, accountability and bias makes them particularly
relevant for a system that executes autonomous actions on a user's machine, handles
sensitive screen data, and runs entirely on local hardware. Unlike other more abstract
ethical guidelines, Google's Principles were designed with real product development
in mind and address the complete lifecycle of an AI system, making them the most
suitable framework for the ethical analysis of a functional system like this one.

Furthermore, given that this project is built around agentic systems, the analysis
of the safety principle is complemented by Google's Secure AI Framework (SAIF)
\cite{saif}, a framework specifically designed to address the security risks of AI
systems, while its objective is very broad, it targets agentic systems too. The SAIF
identifies two primary risks in this context, Sensitive Data Disclosure and Rogue
Actions.

\subsection{Be Socially Beneficial}

AI systems are to provide a tangible value to the users and society, while still
minimising potential harms. This goes beyond functionality, including who is the
end beneficiary, who is excluded from the target audience and what is the broader
impact of the system. In this project, this principle translates into evaluating
what societal problem this system aims to solve and how effectively it does so.

\subsection{Avoid Creating or Reinforcing Unfair Bias}

AI systems can introduce bias in many ways, due to the data they were trained on.
In the context of autonomous agents, bias can lead to an inconsistent behaviour,
a misinterpretation of the user intent or to make certain actions that do not have
the user as the final beneficiary, lacking transparency in the system. Here for
this principle, it should be evaluated the steps that have been taken to ensure
the system is fair and inclusive.

\subsection{Be Built and Tested for Safety}

Safety in AI systems refer to the ability to operate accordingly to the user intent
without causing any harm or disruption, especially in systems that interact with
the real world. In desktop automation, this is crucial, where a misclick can result
in irreversible changes in the user's environment. This principle translates into
evaluating how safety has been prioritised in the design and development of the
system.

\subsection{Be Accountable to People}

In this principle, accountability refers the ability to understand, audit and be
responsible for the decisions made by the AI system. This in autonomous agents is
especially important, where the system can act on behalf the user without necessarily
explaining all the step-by-step reasoning. In this project, this principle translates
into evaluating how the system provides transparency and user control.

\subsection{Incorporate Privacy Design Principles}

The term privacy by design refers to ensure that privacy should be built in the
system from the ground up, rather than being added afterwards once is completed.
In the context of AI systems that interact with the user's desktop, this is
particularly relevant since the system has complete access to sensitive information,
such as personal files, and the content visible in the screen. Therefore, the key
answer to this principle is to reflect on how this system protects user data and
privacy.

\subsection{Uphold High Standards of Scientific Excellence}

When developing systems in AI research, it is important that they comply with
privacy first and security first rules, however, it is important that these systems
are grounded in rigorous research, empirical evaluation and transparent methodology.
This implies that any claim about the system's capabilities should be backed by
evidence and should be justified through scientific reasoning.

\subsection{Be Made Available for Uses that Accord with These Principles}

This principle refers to the responsibility of ensuring that an AI system is only
made available for uses that align with ethical guidelines. This prevents it from
being reused for malicious or unintended intents, therefore the final applications
of this system should be outlined accordingly.

\vspace{3em}
\section{Project Analysis}

The following section evaluates this system against each of the seven Google AI Principles. For each principle, a set of metrics is defined and scored on a scale from 1 to 5, where this applies:

\begin{table}[h]
\centering
\begin{tabular}{|c|l|}
\hline
\textbf{Score} & \textbf{Definition} \\
\hline
1 & The dimension is barely considered, with no concrete implementation \\
2 & The dimension is partially addressed but with significant limitations \\
3 & The dimension is addressed but with some notable gaps or limitations \\
4 & The dimension is well addressed with only minor limitations \\
5 & The dimension is fully and satisfactorily addressed \\
\hline
\end{tabular}
\caption{Scoring rubric for the ethical dimension evaluation.}
\end{table}

\subsection{Be Socially Beneficial}

This principle evaluates the tangible value that this system provides to users and society. The following dimensions have been considered:

\begin{itemize}
    \item \textbf{Accessibility for motor-disabled users:} This system addresses a concrete and underserved accessibility problem, users with motor disabilities who cannot operate a traditional mouse or keyboard have very limited options for autonomous desktop control. This system shows that a fully local AI system is achievable on consumer hardware, which has a clear social value regardless of its research stage, for this reason, a score of 4/5 is assigned.

    \item \textbf{Hardware accessibility:} The entire pipeline fits within 8 GB of VRAM, making it viable on consumer hardware machines rather than requiring expensive infrastructure. This design decision ensures that the system is not restricted to users with cutting edge hardware, therefore a score of 5/5 is assigned.

    \item \textbf{Open-source research contribution:} It is important to note that this is a research project that does not include a graphical user interface, which limits its immediate accessibility to more experienced users. However, as an open-source academic project, it lays the technical foundation for future work that might make this kind of tool available for far more users, contributing on accessible and private AI systems. Despite this, given this limitation, a score of 3/5 is assigned.
\end{itemize}

\pagebreak
\subsection{Avoid Creating or Reinforcing Unfair Bias}

This principle evaluates the steps taken to ensure the system is fair and inclusive, without any predefined bias. The following dimensions have been considered:

\begin{itemize}
    \item \textbf{Determinism of the planning layer (PDDL):} Bias is directly addressed with the PDDL planning agent. Unlike purely LLM-based systems, where the model response can vary depending on context or biases in the training data, a PDDL plan is deterministic, given the same intent and arguments, it will always produce the same sequence of actions. For this reason, a score of 5/5 is assigned.

    \item \textbf{Independence from user-specific training data:} This system does not rely on any user-specific data during inference, meaning that no personal information is used to influence the decisions taken by the pipeline. The underlying models were trained on publicly available data, and no fine-tuning on user data is performed at any stage. However, there might be some hereditary bias from the Qwen models, although is insignificant, given this, a score of 4/5 is assigned.

    \item \textbf{Prompt design for bias mitigation:} All prompts across the pipeline follow a consistent structure: they have a maximum of 350 characters, contain no personal information, and always lead with a clear objective followed by a format example, synonyms or descriptive definitions of the elements involved. This reduces the likelihood of the model interpreting instructions in unintended or biased ways, therefore a score of 4/5 is assigned.
\end{itemize}

\subsection{Be Built and Tested for Safety}

This principle evaluates how safety has been prioritised in the design and development of this system. The following dimensions have been considered:

\begin{itemize}
    \item \textbf{Sensitive Data Disclosure risk mitigation (SAIF):} Google's Secure AI Framework (SAIF) identifies Sensitive Data Disclosure as one of the most important risks in agentic systems, referring to the exposure of private or confidential data through the agent's privileged access. In this system, everything runs locally with no external APIs. Therefore, a score of 5/5 is assigned.

    \item \textbf{Rogue Actions risk mitigation (SAIF):} The SAIF also identifies Rogue Actions as a critical risk, referring to unintended actions executed by an agent, whether accidental or malicious. The controlled action space enforced by the PDDL templates ensures that this system cannot execute actions outside of its predefined scope, and the QA Agent provides an additional verification step to detect any misalignment between the intended and the actual outcome. For this reason, a score of 4/5 is assigned.

    \item \textbf{Controlled action space (PDDL templates):} The use of predefined PDDL task templates means this system can only execute a controlled and well-defined set of actions, preventing it from taking arbitrary or unexpected steps outside of its intended scope. This is a deliberate architectural decision that significantly reduces the risk of unintended or harmful behaviour, therefore a score of 5/5 is assigned.

    \item \textbf{Least-privilege principle per agent:} Each agent in the pipeline only has access to what is strictly necessary for its specific task. This is why, there are 8 agents in the pipeline, despite being a high number, this was thought from the beginning of the project, before the development of the system, to ensure that no agent could do more than the necessary actions. Therefore, a 5/5 is assigned, due to the security by design implementation.
\end{itemize}

\subsection{Be Accountable to People}

This principle evaluates how this system provides transparency and user control. The following dimensions have been considered:

\begin{itemize}
    \item \textbf{Auditability of the PDDL plan:} The PDDL plan generated for each task is fully auditable, meaning that every step the system intends to execute is formally defined and can be inspected. For this reason, a score of 5/5 is assigned.

    \item \textbf{Lack of pre-execution user confirmation:} It is worth acknowledging that this system currently executes plans directly without requesting explicit user confirmation beforehand, which is a limitation from an accountability perspective. Given this limitation, a score of 2/5 is assigned.
\end{itemize}

\subsection{Incorporate Privacy Design Principles}

This principle evaluates how this system protects user data and privacy. The following dimensions have been considered:

\begin{itemize}
    \item \textbf{Full local execution:} All models run locally via Ollama, no screenshots or screen data are ever sent to external servers nor kept in device memory, and the system requires no internet connection to operate. This was thought with the privacy by design concept in mind, giving this dimension a 5/5.

    \item \textbf{Open-source models:} The models used across the pipeline belong to the Qwen3 family, which are fully open source. This means their architecture, training process and behaviour are publicly auditable. This reinforces the privacy guarantees of the system, for this reason, a score of 5/5 is assigned.
\end{itemize}

\subsection{Uphold High Standards of Scientific Excellence}

This principle evaluates whether this system is grounded in rigorous research and transparent methodology. The following dimensions have been considered:

\begin{itemize}
    \item \textbf{Grounding in planning and VLM literature:} The integration of formal PDDL planning with a multi-agent pipeline, ensure that the most optimal path is chosen every time, making choices that in some cases might be better than other LLM based systems. Moreover, this gives backed evidence on the systems capabilities, however, the system is limited to the templates chosen, therefore a score of 4/5 is assigned.

    \item \textbf{Empirical benchmarking against competitors:} This system is being evaluated against state-of-the-art cloud and local competitors, reducing the potential bias in own comparison statistics. Strengthening the scientific validity of the evaluation. Therefore a score of 5/5 is assigned.
\end{itemize}

\subsection{Be Made Available for Uses that Accord with These Principles}

This principle evaluates whether this system is only made available for uses that align with ethical guidelines. The following dimension has been considered:

\begin{itemize}
    \item \textbf{Restricted action space by design:} The architecture itself acts as a natural constraint, as it only supports a predefined set of task templates, meaning it cannot be repurposed for any harmful automation. Moreover, the minimum privilege access of agents, ensure that in case of any misuse, it is contained as much as possible. Although, the PDDL templates where not thought as a security measure, but more as limitation of the scope of the project, therefore a score of 4/5 is assigned.
\end{itemize}

\section{Impact Scale of Ethical Aspects}

The following table summarises the scores obtained for each of the seven Google AI Principles, based on the dimensions evaluated in the previous section. The overall score is calculated as the average of all principle averages.

\begin{table}[h]
\centering
\begin{tabular}{|l|c|}
\hline
\textbf{Google AI Principle} & \textbf{Avg. Score} \\
\hline
Be Socially Beneficial & 4.0 / 5 \\
Avoid Creating or Reinforcing Unfair Bias & 4.3 / 5 \\
Be Built and Tested for Safety & 4.75 / 5 \\
Be Accountable to People & 3.5 / 5 \\
Incorporate Privacy Design Principles & 5.0 / 5 \\
Uphold High Standards of Scientific Excellence & 4.5 / 5 \\
Be Made Available for Uses that Accord with These Principles & 4.0 / 5 \\
\hline
\textbf{Overall} & \textbf{4.29 / 5} \\
\hline
\end{tabular}
\caption{Summary of Google AI Principles scores.}
\end{table}

The results show that this system performs particularly well in privacy, safety and availability, all of which are directly enforced by its local-first architecture. The lowest score corresponds to accountability, which is an acknowledged limitation of the implementation given the absence of a pre-execution confirmation step. Overall, with a score of 4.29 out of 5, this system demonstrates a strong alignment with Google's AI Principles, with clear areas for improvement identified for future iterations.